% Source: https://artofproblemsolving.com/community/c3h1249717_a_hard_problem_about_arithmetic
% Stable identity: aops-topic-1249717-q01
If $P$ has $k$ divisors, pairing each divisor $d$ with $\dfrac {P}{d}$ shows that their product is $P^{\dfrac {k}{2}}$. So $P^{\dfrac {k}{2}}=P^2$; $P=1$ fails the sum condition, so $k=4$, i.e. $P=p^3$ or $P=pq$ with distinct primes $p<q$.

If $P=p^3$: $1+p+p^2+p^3=2p^3$ gives $1+p+p^2=p^3$, impossible modulo $p$.

If $P=pq$: $(1+p)(1+q)=2pq$, i.e. $pq-p-q=1$, $(p-1)(q-1)=2$, so $p=2$, $q=3$.

The only such number is $\boxed{6}$.
